% ---------------------------------------------------------------------------
% Chapter 9: The First Checkpoint: Level-1 Trigger
% Generated from the outline; edit freely, this file is now the source.
% ---------------------------------------------------------------------------
\chapter{The First Checkpoint: Level-1 Trigger}
\chaptermark{The First Checkpoint}
\label{ch:first_checkpoint}

\begin{journeybox}
At the first checkpoint a guard with a coarse ledger has a few microseconds to
decide whether the family may continue at all. He sees only towers, not people,
and his thresholds are the first place where families are lost with a
probability rather than a certainty.
\end{journeybox}

\physicsline{calorimeter trigger towers, Layer-1 and Layer-2, $9\times9$ tower jets, L1 pileup subtraction, uGT menu, efficiency turn-ons}

\section{Forty million times per second: the L1 problem}
\label{sec:first_checkpoint:forty_million_times_per_second_the_l1_pr}
\todo{Write this section.}

\section{Trigger towers and the L1 jet algorithm}
\label{sec:first_checkpoint:trigger_towers_and_the_l1_jet_algorithm}
\todo{Write this section.}

\section{Level-1 pileup subtraction and calibration look-up tables}
\label{sec:first_checkpoint:level_1_pileup_subtraction_and_calibrati}
\todo{Write this section.}

\section{Efficiency turn-on curves: the mathematics of a threshold}
\label{sec:first_checkpoint:efficiency_turn_on_curves_the_mathematic}
\todo{Write this section.}

\section{Our jet at this stage: did our event pass, and on which L1 seed?}
\label{sec:first_checkpoint:our_jet_at_this_stage_did_our_event_pass}
\todo{Numbers, plots and event display of the $b$-jet at this stage.}

\begin{ledger}{The First Checkpoint}
  \ledgerrow{before this stage}{}{}{}{--}
  \ledgerrow{after this stage}{}{}{}{}
\end{ledger}

\section{Further reading}
\label{sec:first_checkpoint:further_reading}
Official and theory references for this stage: \cite{CMS:2020cmk}, \cite{CMS:2016ngn}, \cite{CMS:2020jkc}.

\begin{pitfalls}
\todo{Pitfalls and FAQs for newcomers at this stage.}
\end{pitfalls}

\begin{exercises}
\item \todo{Exercise 1.}
\item \todo{Exercise 2.}
\end{exercises}
